% 水星（综述）
% license CCBYSA3
% type Wiki

本文根据 CC-BY-SA 协议转载翻译自维基百科\href{https://en.wikipedia.org/wiki/Mercury_(planet)}{相关文章}。

\begin{figure}[ht]
\centering
\includegraphics[width=8cm]{./figures/8ad510c3701a820c.png}
\caption{} \label{fig_Mercur_1}
\end{figure}
Mercury 是距离太阳最近的行星，也是太阳系中最小的行星。它是一颗岩质行星，具有极其稀薄的大气层，其表面重力略高于火星。水星的表面与地球的月球相似，布满大量撞击坑，并具有由逆断层形成的广阔断崖系统（rupes），以及由抛射物形成的明亮光条结构。其最大撞击盆地——卡洛里斯平原（Caloris Planitia）——直径为 \(1{,}550\,\mathrm{km}\)（\(960\,\mathrm{mi}\)），约为整颗行星直径 \(4{,}880\,\mathrm{km}\)（\(3{,}030\,\mathrm{mi}\)）的三分之一。作为最靠内侧运行的行星，它在地球的天空中总是接近太阳出现，呈现为“晨星”或“昏星”。它也是从地球前往所需 \(\Delta v\) 最大的行星，同时亦是往返太阳系中其他行星时所需能量最高者。

水星的恒星年（\(88.0\) 个地球日）与恒星日（\(58.65\) 个地球日）构成 \(3:2\) 的自转—公转共振。因此，一个太阳日（从日出到下一次日出）约持续 \(176\) 个地球日，即其恒星年的两倍。这意味着水星某一侧将在一个水星年（\(88\) 个地球日）内持续处于日照之下；而在接下来的轨道周期中的另一个 \(88\) 个地球日内，该侧将一直处于黑暗之中，直到下一次日出。行星表面之上存在极度稀薄的外逸层以及微弱但足以偏转太阳风的磁场。由于水星轨道偏心率极高，其表面所接受的日照强度与温度变化极为剧烈：赤道地区夜间温度可达 \(-170^{\circ}\mathrm{C}\)（\(-270^{\circ}\mathrm{F}\)），而白昼可升至 \(420^{\circ}\mathrm{C}\)（\(790^{\circ}\mathrm{F}\)）。由于轴倾角极小，水星两极区域永久处于阴影之中，这强烈暗示其撞击坑内可能存在水冰。

与太阳系其他行星相同，水星形成于约 \(4.5\times10^{9}\) 年前。关于其起源与演化存在诸多相互竞争的假说，其中一些涉及与微行星碰撞及岩石蒸发过程；截至 2020 年代初，水星地质历史的许多宏观细节仍在研究中，或需等待更多深空探测数据。其地幔成分高度均一，这表明水星早期可能存在类似月球的岩浆海。根据当前模型，水星可能具有固态硅酸盐地壳与地幔，其下为固态外核、更深处的液态内核层，以及固态内核。约在七至八十亿年后，随着太阳演化为红巨星，水星预计将被毁灭，金星亦然，地球与月球也可能遭遇同样命运\(^\text{[20]}\)。

水星是一颗“古典行星”，自古以来便被人类持续观测并识别为行星（或“游星”）。英文名称 Mercury 源自古罗马神祇 Mercurius，即商业与沟通之神，也是众神的信使。1974 年，“水手 10 号”（Mariner 10）完成了首次成功飞掠水星的任务，此后 MESSENGER 与 BepiColombo 轨道器也相继对其进行访问与探测。
\subsection{命名}

在历史上，人类会根据水星作为“昏星”或“晨星”出现的不同情况而称呼它不同的名字。到公元前约 350 年，古希腊人已经意识到这两颗“星”实际上是同一个天体[21]。他们称这颗行星为 Στίλβων（Stilbōn），意为“闪烁者”，以及 Ἑρμής（Hermēs），因其迅速而短暂的移动\,[22]；这一名称在现代希腊语中依然保留（Ερμής Ermis）[23]。罗马人则以他们的信使之神墨丘利（Mercury，拉丁语 Mercurius）为其命名，因为水星在天空中的移动速度比任何其他行星都快[21][24]。不过，根据老普林尼的记载，也有人将这颗行星与阿波罗相联系[25]。

水星的天文学符号是赫尔墨斯权杖（caduceus）的风格化形式；在 16 世纪，该符号加入了一个基督教十字，其形状为：☿[26][27]
\subsection{物理特性}
\begin{figure}[ht]
\centering
\includegraphics[width=8cm]{./figures/23fb842f59355847.png}
\caption{按比例绘制的水星及内太阳系具有行星质量的天体。从左至右依次为：水星、金星、地球、月球、火星和谷神星。} \label{fig_Mercur_2}
\end{figure}
水星是太阳系中四颗类地行星之一，这意味着它与地球一样，是一颗岩质天体。它是太阳系中最小的行星，赤道半径为 \(2{,}439.7\,\mathrm{km}\)（\(1{,}516.0\,\mathrm{mi}\)）[4]。水星的体积也小于——但质量大于——太阳系中最大的两个天然卫星：木卫三（盖尼米得）和土卫六（泰坦）。水星的组成大约为 70\% 的金属物质与 30\% 的硅酸盐物质[28]。
\subsubsection{内部结构}
\begin{figure}[ht]
\centering
\includegraphics[width=8cm]{./figures/697a86a5c3dcc15c.png}
\caption{水星的内部结构与磁场} \label{fig_Mercur_3}
\end{figure}
水星似乎具有一层固态硅酸盐地壳与地幔，它们覆盖着固态金属外核层、更深的液态内核层，以及固态内核[29][30]。富铁内核的成分仍不确定，但很可能包含镍、硅，或许还有硫和碳，并含有微量其他元素[31]。水星的密度是太阳系中第二高的，为 \(5.427\,\mathrm{g/cm^{3}}\)，仅略低于地球的 \(5.515\,\mathrm{g/cm^{3}}\)[4]。如果将两颗行星的重力压缩效应排除，水星的构成物质将比地球更致密：其未压缩密度约为\(5.3\,\mathrm{g/cm^{3}}\)，而地球为 \(4.4\,\mathrm{g/cm^{3}}\)[32]。水星的密度可用于推断其内部结构。尽管地球的高密度在很大程度上源于重力压缩（尤其在核心处），但水星要小得多，其内部区域并未受到同等程度的压缩。因此，为了呈现如此高的密度，其核心必须既巨大又富含铁[33]。

基于满足惯性矩因子 \(0.346\pm0.014\) 的内部结构模型，水星的核心半径估计为 \(2{,}020\pm30\,\mathrm{km}\)（\(1{,}255\pm19\,\mathrm{mi}\)）[9][34]。因此，水星核心占其总体积约 57\%；相比之下，地球仅为 17\%。2007 年发表的研究表明水星具有熔融内核\,[35][36]。地幔—地壳层总厚度约为 \(420\,\mathrm{km}\)（\(260\,\mathrm{mi}\)）\,[37]。关于地壳厚度，各研究推算存在差异：Mariner~10 与 MESSENGER 探测数据给出的估计为 \(35\,\mathrm{km}\)（\(22\,\mathrm{mi}\)），而基于 Airy 均衡假说的模型则给出 \(26\pm11\,\mathrm{km}\)（\(16.2\pm6.8\,\mathrm{mi}\)）[38][39][40]。水星表面的一个显著特征是大量狭长的脊状地形，可延伸数百公里。一般认为这些构造形成于水星核心与地幔冷却并收缩的时期，而那时地壳已处于固态[41][42][43]。

水星的核心含铁量在太阳系行星中最高，为解释这一特征已提出多种理论。最广为接受的观点认为，水星最初的金属—硅酸盐比例与普通球粒陨石相似，而后者被认为代表太阳系岩质物质的典型组成；当时水星的质量大约为其当前质量的 2.25 倍[44]。在太阳系早期，水星可能被一颗质量约为其六分之一、直径达数千公里的微行星撞击[44]。这一撞击会剥离其大部分原始地壳与地幔，使核心成为主要残留成分\,[44]。一种类似的机制，即巨撞假说（giant impact hypothesis），也被用于解释地球月球的形成[44]。

另一种假说认为，水星可能形成于太阳星云的早期阶段，那时太阳的能量输出尚未稳定。其初始质量可能是当前质量的两倍，但随着原太阳（protosun）收缩，水星附近的温度可能达到 \(2{,}500\sim3{,}500\,\mathrm{K}\)，甚至高达 \(10{,}000\,\mathrm{K}\)[45]。在如此温度下，水星的大部分表面岩石会被汽化，形成“岩石蒸汽”大气，并可能被太阳风带走\,[45]。第三种假说则认为，太阳星云对水星聚积的固体颗粒施加阻力，导致较轻的颗粒从聚积物质中被移除，从而未被水星捕获[46]。

每一种假说均预测不同的表面化学组成，因此已有两项深空任务专门针对这一问题进行观测。第一项任务 MESSENGER（2015 年结束）发现水星表面的钾与硫含量高于预期，这表明巨撞假说与地壳—地幔蒸发模型可能并未发生，因为这些事件的极端高温会使钾与硫完全挥发[47]。BepiColombo 将于 2025 年抵达水星，并将继续对这些假说进行检验[48]。迄今为止的结果似乎支持第三种假说；然而，仍需进一步分析数据[49]。
\subsubsection{表面地质}
\begin{figure}[ht]
\centering
\includegraphics[width=8cm]{./figures/02332822930c85ad.png}
\caption{利用 NASA 和 USGS 数据在 Blender 中生成的水星渲染图} \label{fig_Mercur_4}
\end{figure}
水星的表面在外观上与月球相似，呈现广阔的类月海平原以及大量撞击坑，这表明其地质活动已停滞达数十亿年之久。与火星或月球相比，水星表面更加不均质；火星与月球都具有大面积的相似地质单元，例如月海与高原\,[50]。反照率特征是指具有显著不同反射率的区域，其中包括撞击坑、由此形成的抛射物以及光条系统。较大的反照率特征通常对应反射率更高的平原区域\,[51]。水星具有“皱褶脊”（dorsa）、类似月球的高地、山脉（montes）、平原（planitiae）、断崖（rupes）以及峡谷（valles）[52][53]。
\begin{figure}[ht]
\centering
\includegraphics[width=8cm]{./figures/8c611b128a504de7.png}
\caption{MESSENGER 对水星表面进行的 MASCS 光谱扫描} \label{fig_Mercur_5}
\end{figure}
水星的地幔在化学成分上具有不均一性，这表明该行星在早期历史中经历过岩浆海阶段。矿物的结晶与对流翻转导致形成了一层具有化学分层、成分不均的地壳，而这种大尺度的化学组成差异也可在其表面观测到。地壳含铁量低但含硫量高，这源于水星早期比其他类地行星更强的化学还原环境。水星表面主要由贫铁的辉石与橄榄石组成，分别以顽辉石（enstatite）与镁橄榄石（forsterite）为代表，同时还包含富钠的斜长石，以及由镁、钙与硫化铁混合构成的矿物。地壳中反照率较低的区域含有高含量的碳，最可能以石墨的形式存在[54][55]。

水星地表特征的命名来自多种来源，并依据 IAU 行星命名系统设定。以人物命名的仅限于已故者。撞击坑以在艺术、音乐、绘画与文学等领域做出杰出或基础性贡献的人物命名。脊或皱褶（dorsa）以对水星研究有贡献的科学家命名。洼地或槽谷（fossae）以建筑作品命名。山脉（montes）以各语言中表示“热”的词语命名。平原（planitiae）以各语言中“水星”的名称命名。断崖（rupēs）以科学考察的船只命名。峡谷（valles）以古代被废弃的城市、城镇或聚落命名[56]。

\textbf{撞击盆地与撞击坑}
\begin{figure}[ht]
\centering
\includegraphics[width=8cm]{./figures/cc1a7ee3ff642536.png}
\caption{} \label{fig_Mercur_6}
\end{figure}













